\documentclass[10pt,twocolumn]{article}
\usepackage[T1]{fontenc}
\usepackage[utf8]{inputenc}
\usepackage{lmodern}
\usepackage[margin=1.8cm,top=2.4cm,bottom=2cm,columnsep=0.6cm]{geometry}
\usepackage{fancyhdr}
\usepackage{titlesec}
\usepackage{booktabs}
\usepackage{amsmath,amssymb,amsfonts}
\usepackage{microtype}
\usepackage{xcolor}
\usepackage{mdframed}
\usepackage{parskip}
\usepackage{enumitem}
\usepackage{hyperref}

\definecolor{rulered}{RGB}{180,30,30}
\definecolor{boxgray}{RGB}{245,245,245}
\definecolor{keywordgray}{RGB}{100,100,100}

\pagestyle{fancy}
\fancyhf{}
\fancyhead[L]{\textsc{\small Claude Mag}}
\fancyhead[C]{\small\textit{Machine Learning Research Digest}}
\fancyhead[R]{\small 2026-10-28}
\fancyfoot[C]{\small\thepage}
\renewcommand{\headrulewidth}{0.8pt}

\titleformat{\section}{\normalsize\bfseries\color{rulered}}{}{0pt}{}%
  [\vspace{-4pt}\color{rulered}\rule{\columnwidth}{0.4pt}]
\titlespacing{\section}{0pt}{8pt}{4pt}

\hypersetup{colorlinks=true,urlcolor=rulered,linkcolor=black}

\begin{document}
\setlength{\parindent}{0pt}
\setlength{\parskip}{4pt}

{\small\color{keywordgray}\textbf{Keywords:}
  test-time adaptation \textbullet{}
  robot locomotion \textbullet{}
  reinforcement learning \textbullet{}
  distribution shift}\par\vspace{4pt}

{\LARGE\bfseries\raggedright
  TEMPO: Efficient Test-Time Adaptation For
  Robot Policies}\par\vspace{4pt}

{\large\itshape\raggedright
  Critic Entropy Regularization and Score-Based Memory Enable
  On-Device Policy Adaptation---Up to 80\% Performance Recovery
  on Physical Quadrupeds Under Load}\par\vspace{6pt}

{\small\textbf{By} Motasem Alfarra, Pietro Mazzaglia, Markus Peschl,
  Daniel Dijkman, Christos Louizos}\par\vspace{2pt}

{\small\color{keywordgray}arXiv:2610.20318 \textbullet{}
  NeurIPS 2026}
\par\vspace{2pt}\noindent{\color{rulered}\rule{\columnwidth}{0.6pt}}\vspace{6pt}

Reinforcement-learned locomotion policies trained in simulation
degrade in deployment when conditions shift---payload changes the
center of mass, terrain modifies effective friction, sensors drift
with temperature. TEMPO addresses this with a surprisingly effective
combination: an entropy regularizer on the critic that prevents
value collapse under distribution shift, paired with a score-based
memory buffer that upweights hard observations during on-device
adaptation. Across three quadruped platforms (Go1, Go2, G1) and
four shift types in MuJoCo Playground, TEMPO yields up to 80\%
improvement in the single-device setting; on a physical Unitree Go2
carrying an added payload, TEMPO recovers pre-shift performance
within minutes.

\begin{mdframed}[backgroundcolor=boxgray,linecolor=rulered,linewidth=1pt,
    innerleftmargin=6pt,innerrightmargin=6pt,
    innertopmargin=6pt,innerbottommargin=6pt]
\textbf{\small AT A GLANCE}\par\vspace{4pt}
\begin{description}[leftmargin=0pt,labelsep=4pt,itemsep=2pt,parsep=0pt,
    font=\small\bfseries]
  \item[Problem:] Sim-to-real locomotion policies degrade when
    deployment conditions (payload, terrain, sensors) differ from
    training---and no simulator is available at test time.
  \item[Why it matters:] Real-world deployment requires policies that
    adapt autonomously; simulator-based fine-tuning is impractical
    on hardware.
  \item[Method:] Critic entropy regularizer prevents value collapse;
    score-based memory oversamples observations the policy finds hard.
  \item[Result:] Up to 80\% performance improvement in simulation;
    physical Unitree Go2 recovers from payload shift.
  \item[Evidence:] Go1, Go2, G1 in MuJoCo Playground under four
    distribution shifts; physical Unitree Go2 experiment.
\end{description}
\end{mdframed}

\section{The Big Picture}

Reinforcement learning in simulation remains the dominant approach for
robot locomotion, but the sim-to-real gap means deployed policies face
conditions their training never covered. The standard remedy---domain
randomisation---widens the training distribution but cannot anticipate
all deployment shifts. Once deployed, the robot's own experience is the
only data available for further improvement.

Test-time adaptation (TTA) of neural networks has been studied for
supervised classification, but RL policies present a harder challenge:
there are no ground-truth labels, the reward is expensive to evaluate,
the critic can diverge under distribution shift, and compute is
constrained to a single embedded device. TEMPO is the first method to
address all four constraints simultaneously.

\section{Why Existing Approaches Fall Short}

\textbf{Standard PPO fine-tuning at test time:}
Without a simulator, PPO must use on-device rollouts directly. The
off-policy nature of the limited experience combined with a shifting
observation distribution causes the critic to diverge, destabilising
policy updates.

\textbf{Domain randomisation:}
Broad randomisation reduces peak sim performance and still fails for
shifts outside the randomised range (e.g., a specific payload not
anticipated during training).

\textbf{Meta-learning (MAML, PEARL):}
Meta-learning requires a distribution over adaptation tasks at training
time---unavailable when the real-world distribution is unknown---and
inner-loop gradient computation that is expensive on embedded hardware.

\section{Core Mechanism}

TEMPO adds two components on top of PPO, requiring no changes to the
policy network or training procedure:

\textbf{1. Critic Entropy Regularizer.}
The critic loss is augmented with an entropy penalty:
\[
\mathcal{L}_V = \mathcal{L}_\mathrm{TD} + \beta\, H(V_\phi(s))
\]
where $H(V_\phi(s))$ is the entropy of the value distribution across
a replay batch. This prevents the critic from collapsing to a
near-constant value function under distribution shift, which would
zero out policy gradient signals.

\textbf{2. Score-Based Memory.}
Each observation $s$ in the replay buffer is assigned a score
proportional to its TD error:
\[
\mathrm{score}(s) \propto |r + \gamma V_\phi(s') - V_\phi(s)|
\]
Observations with high scores---those where the critic is most
wrong---are sampled more frequently, biasing updates toward the
hardest parts of the new distribution. Scores are updated online as
the critic improves.

\section{Technical Formulation}

The full TEMPO objective at each adaptation step is:
\[
\mathcal{L}_\mathrm{TEMPO} = \mathcal{L}_\mathrm{PPO} +
  \lambda_V \bigl[\mathcal{L}_\mathrm{TD} + \beta\, H(V_\phi)\bigr]
\]
where $\mathcal{L}_\mathrm{PPO}$ is the clipped surrogate objective
and $\lambda_V$ balances the critic and actor losses.

The score-based sampler maintains a priority queue over $N$ stored
transitions, replacing the oldest with probability proportional to
inverse score (keeping recent hard examples longer). Sampling
probability for transition $i$ is:
\[
p_i = \frac{\mathrm{score}(s_i)^\alpha}
  {\sum_j \mathrm{score}(s_j)^\alpha}, \quad \alpha = 0.6
\]
(Priority exponent $\alpha$ follows the PER convention.)

\section{Learning or Inference Procedure}

\textbf{Offline training (simulation):} Standard PPO with domain
randomisation. No modifications to the training procedure.

\textbf{Online adaptation (deployment):} At each deployment step:
\begin{enumerate}[itemsep=2pt,parsep=0pt]
  \item Execute the current policy on the robot; collect transitions.
  \item Update transition scores from critic TD errors.
  \item Sample a batch from the priority queue.
  \item Apply one PPO + entropy-regularised critic update step.
\end{enumerate}
Total compute per step: equivalent to one standard PPO update,
feasible on the onboard CPU/GPU of a Unitree robot.

\section{What the Guarantee Says}

Under mild assumptions (bounded reward, Lipschitz value function,
score ergodicity), the entropy regularizer prevents the critic from
collapsing to a degenerate fixed point during distribution shift:
\[
H(V_\phi(s)) \geq H_\mathrm{min} > 0 \;\;\text{for all } t
\]
ensuring policy gradient variance stays non-zero. This provides a
necessary condition for adaptation progress; a convergence rate
depends on the magnitude of the distribution shift and the data
collection speed.

\section{Experimental Findings}

\begin{itemize}[itemsep=2pt,parsep=0pt]
  \item \textbf{Go2 (MuJoCo Playground, payload shift):} TEMPO reaches
    80\% of pre-shift performance within 5 minutes of on-device
    interaction; standard PPO fine-tuning diverges.
  \item \textbf{Go1, G1 (MuJoCo Playground, four shift types):}
    Consistent gains across all platforms and shift types; average
    $+42\%$ over non-adapted baseline.
  \item \textbf{Physical Unitree Go2:} TEMPO with a 2\,kg added
    payload recovers stable locomotion within 10 minutes; the
    non-adapted policy fails to sustain a trot gait.
  \item \textbf{Baseline comparison:} Outperforms domain randomisation,
    meta-RL (MAML), and test-time normalisation across all conditions.
\end{itemize}

\section{Ablations and Interpretation}

\textbf{Entropy regularizer only:}
Without score-based memory, the critic stabilises but adaptation is
slower ($-25\%$ performance at 5\,min); both components are needed.

\textbf{Score-based memory only:}
Without entropy regularization, the critic still diverges in 3 of 4
shift scenarios, confirming that stability must come before
curriculum selection.

\textbf{Priority exponent $\alpha$:}
Performance is robust to $\alpha \in [0.4, 0.7]$; $\alpha < 0.4$
approaches uniform sampling and loses most gains.

\textbf{Batch size:}
Adaptation improves up to batch size 64 and plateaus, consistent
with the variance reduction argument underlying critic entropy
regularization.

\section*{Reference}

Motasem Alfarra, Pietro Mazzaglia, Markus Peschl, Daniel Dijkman,
Christos Louizos.
\textbf{Efficient Test-Time Adaptation For Robot Policies}.
arXiv:2610.20318. NeurIPS 2026.
\url{https://arxiv.org/abs/2610.20318}

\end{document}
